\documentclass[11pt]{amsart}
\usepackage{amsmath,amssymb,amsthm,booktabs}
\usepackage[margin=1.15in]{geometry}
\usepackage[hidelinks]{hyperref}
\newtheorem{theorem}{Theorem}[section]
\newtheorem{proposition}[theorem]{Proposition}
\newtheorem{lemma}[theorem]{Lemma}
\newtheorem{corollary}[theorem]{Corollary}
\theoremstyle{definition}
\newtheorem{definition}[theorem]{Definition}
\theoremstyle{remark}
\newtheorem{remark}[theorem]{Remark}
\newcommand{\R}{\mathbb{R}}
\newcommand{\one}{\mathbf{1}}
\newcommand{\doi}[1]{\href{https://doi.org/#1}{doi:\nolinkurl{#1}}}
\setlength{\emergencystretch}{2.5em}

\title[The regularity threshold of positive quadratic stochastic operators]{On the regularity threshold of uniformly positive quadratic stochastic operators}
\author{Yitzchak Shmalo}
\address{Einstein Institute of Mathematics, The Hebrew University of Jerusalem, Jerusalem, Israel}
\email{yspennstate@gmail.com}
\date{November 2026}
\subjclass[2020]{37N25, 47H10, 92D25}
\keywords{quadratic stochastic operator, regularity, heredity coefficients, Fisher--Rao metric, contraction, periodic orbit}

\raggedbottom
\begin{document}
\begin{abstract}
A quadratic stochastic operator on the simplex of probability vectors in $\R^m$ is regular if all of its trajectories
converge to one point. Ganikhodzhaev and Sarymsakov showed that the operator is regular when all of its heredity
coefficients exceed $1/(2m)$, and Ganikhodzhaev and Rozikov asked for the least number $\lambda_m$ with this property.
For $m=2$ we prove $\lambda_2=(3-\sqrt7)/4$. For every $m\ge3$ we prove that $\lambda_m$ lies between
$(1-1/\sqrt{m+1})/(2m)$ and $2/(3m+2+\sqrt{m^2+12m+4})$, which is less than $1/(2m)$; in particular
$2m\lambda_m\to1$. The upper bound comes from a contraction of the operator in the Fisher--Rao metric, which
rests on two inequalities: the derivative of a quadratic stochastic operator at most doubles the Fisher information of
a tangent vector, and the uniform-mixture map $q\mapsto\mu\mathbf 1+(1-m\mu)q$ multiplies it by at most
$(1-m\mu)^2/[1-(m-2)\mu]$, where $\mu$ is the least heredity coefficient. The lower bound comes from an explicit family of operators with strictly positive coefficients
and a two-cycle.
\end{abstract}
\maketitle

\section{Introduction}\label{sec:intro}

Let $m\ge2$ and let $S^{m-1}=\{x\in\R^m:x_i\ge0,\ \sum_ix_i=1\}$ be the simplex of probability vectors. A
\emph{quadratic stochastic operator} is a map $V:S^{m-1}\to S^{m-1}$ of the form
\begin{equation}\label{eq:qso}
(Vx)_k=\sum_{i,j=1}^m P_{ij,k}\,x_ix_j,\qquad k=1,\dots,m,
\end{equation}
where the \emph{heredity coefficients} satisfy
\begin{equation}\label{eq:coeff}
P_{ij,k}\ge0,\qquad P_{ij,k}=P_{ji,k},\qquad \sum_{k=1}^mP_{ij,k}=1 .
\end{equation}
In the population model behind this definition, $P_{ij,k}$ is the probability that parents of types $i$ and $j$ have
an offspring of type $k$, and $Vx$ is the distribution of types in the next generation when the present one is $x$; see
the surveys~\cite{GR09,GMR11}. The operator $V$ is \emph{regular} if there is a point $p\in S^{m-1}$ with $V^nx\to p$
for every $x\in S^{m-1}$. The point $p$ is then the only fixed point of $V$. Write
\[
\mu(P)=\min_{i,j,k}P_{ij,k}
\]
for the least heredity coefficient, and note that $\mu(P)\le1/m$ by~\eqref{eq:coeff}.

Ganikhodzhaev and Sarymsakov~\cite{GS88} proved that $V$ is regular when $P_{ij,k}>1/(2m)$ for all $i,j,k$.
Positive quadratic stochastic operators with more than one fixed point are known from the work of Krapivin and
Lyubich~\cite{KL76}, and Fasino and Tudisco~\cite{FT20} recover the bound $1/(2m)$ from the optimal $\ell^1$ Lipschitz
constant of~\eqref{eq:qso}, an ergodicity coefficient in the sense of Dobrushin~\cite{Dobrushin56}. The
survey of Ganikhodzhaev and Rozikov~\cite[Theorem~7 and Problem~12]{GR09}, repeated in~\cite[Theorem~2.28 and
Problem~2.29]{GMR11}, asks for the least number with this property, that is, for
\begin{equation}\label{eq:lambda}
\lambda_m=\inf\bigl\{\lambda\ge0:\ \text{every $V$ with $P_{ij,k}>\lambda$ for all $i,j,k$ is regular}\bigr\}.
\end{equation}
Equivalently, $\lambda_m$ is the supremum of $\mu(P)$ over the non-regular operators, with $\lambda_m=0$ if every
operator with $\mu(P)>0$ is regular. Indeed, if $\lambda$ is admissible in~\eqref{eq:lambda}, then every non-regular
operator has $\mu(P)\le\lambda$; conversely, every number at least the supremum is admissible, since an operator with
all coefficients above it is regular. The same surveys give the value $(3-\sqrt7)/2$ for $m=2$. Saburov and
Saburov~\cite{SS25} restate the question for general operators, where it is open, and answer it for a doubly
stochastic class.

\begin{theorem}\label{thm:binary}
$\lambda_2=(3-\sqrt7)/4=0.0885\dots$. More precisely, every quadratic stochastic operator on $S^1$ with
$\mu(P)\ge(3-\sqrt7)/4$ is regular, and for every $\delta\in\bigl(0,(3-\sqrt7)/4\bigr)$ the operator with heredity
coefficients $P_{11,1}=P_{12,1}=\delta$ and $P_{22,1}=1-\delta$ has $\mu(P)=\delta$ and a two-cycle.
\end{theorem}

The surveys~\cite[Theorem~7(2)]{GR09} and~\cite[Theorem~2.28(ii)]{GMR11} state $\alpha_2=(3-\sqrt7)/2=0.177\dots$
under the same normalization~\eqref{eq:coeff}. By Theorem~\ref{thm:binary} every operator with $\mu(P)\ge(3-\sqrt7)/4$ is
regular, while the family in that theorem has a two-cycle for every $0<\delta<(3-\sqrt7)/4$, so the threshold is
$(3-\sqrt7)/4$. Section~\ref{sec:binary} also shows that the bound of
Theorem~\ref{thm:upper} below, applied with $m=2$, gives $\lambda_2\le(2-\sqrt2)/4=0.146\dots$.

\begin{theorem}\label{thm:upper}
Let $m\ge2$ and
\[
C_m=\frac{2}{3m+2+\sqrt{m^2+12m+4}}=\frac{3m+2-\sqrt{m^2+12m+4}}{4m^2}.
\]
If $\mu(P)>C_m$, then $V$ is a contraction of $S^{m-1}$ in the Fisher--Rao metric, and in particular it is regular.
Hence $\lambda_m\le C_m$, and $C_m<1/(2m)$.
\end{theorem}

\begin{theorem}\label{thm:lower}
Let $m\ge3$ and
\[
L_m=\frac{1-1/\sqrt{m+1}}{2m}.
\]
For every $\ell\in(0,L_m)$ there is a quadratic stochastic operator on $S^{m-1}$ with $\mu(P)=\ell$ and a two-cycle.
Hence $\lambda_m\ge L_m$.
\end{theorem}

\begin{corollary}\label{cor:asym}
For every $m\ge3$, $L_m\le\lambda_m\le C_m<1/(2m)$ (Table~\ref{tab:bounds} gives some values). In particular $2m\lambda_m\to1$ as $m\to\infty$, and
$\lambda_m=1/(2m)-\varepsilon_m$ with $c_1m^{-2}\le\varepsilon_m\le c_2m^{-3/2}$ for absolute constants
$c_1,c_2>0$.
\end{corollary}

\begin{proof}
The first statement combines Theorems~\ref{thm:upper} and~\ref{thm:lower}. As $m\to\infty$,
$1/(2m)-L_m=1/(2m\sqrt{m+1})$, and $3m+2+\sqrt{m^2+12m+4}=4m+8+O(1/m)$, so $1/(2m)-C_m=m^{-2}+O(m^{-3})$.
\end{proof}

\begin{table}[tbp]
\centering
\begin{tabular}{rccc}
\toprule
$m$ & $L_m$ & $C_m$ & $1/(2m)$\\
\midrule
3 & $1/12=0.0833$ & $1/9=0.1111$ & $0.1667$\\
4 & $(5-\sqrt5)/40=0.0691$ & $(7-\sqrt{17})/32=0.0899$ & $0.125$\\
5 & $0.0592$ & $0.0757$ & $0.1$\\
10 & $0.0349$ & $0.0426$ & $0.05$\\
100 & $0.004502$ & $0.004904$ & $0.005$\\
\bottomrule
\end{tabular}
\caption{The bounds of Corollary~\ref{cor:asym} for some $m$.}\label{tab:bounds}
\end{table}

The exact value of $\lambda_m$ for $m\ge3$ remains open. For $m=3,4$, numerical optimization of operators with a
two-cycle or with two distinct fixed points does not go beyond least coefficient $L_m$ (Section~\ref{sec:numerics}).

\subsection*{Related work}
For $m=2$ the trajectories of~\eqref{eq:qso} are those of a quadratic map of $[0,1]$. By Lyubich's analysis of
quadratic interval maps~\cite{Lyubich74}, recalled in~\cite[Theorem~3.1]{KHG22}, this map is regular when the discriminant $D$ of $f(x)-x$ is less
than $4$ and has a two-cycle when $D>4$. The proof of Theorem~\ref{thm:binary} needs, in addition, the largest value
of $D$ under a lower bound on the coefficients and the case $D=4$. Fasino and Tudisco~\cite[Theorem~6.1 and
Corollary~6.2]{FT20} compute the optimal Lipschitz constant of~\eqref{eq:qso} in the $\ell^1$ norm and recover the bound
$1/(2m)$; the same contraction is Proposition~\ref{prop:l1} below. Theorem~\ref{thm:upper} replaces the $\ell^1$ norm
by the Fisher--Rao metric, in which the uniform part of the coefficients improves the constant. Lipschitz estimates for
polynomial operators on the simplex go back to Krapivin and Lyubich~\cite{KL77}, and Saburov~\cite{S26} uses contraction
to prove global asymptotic stability of polynomial stochastic operators.

\subsection*{Organization}
The paper is arranged as follows; the notation is collected in Appendix~\ref{app:notation}.
\begin{center}
\begin{tabular}{lll}
\toprule
Section & Content & Used in\\
\midrule
\ref{sec:prelim} & $V=\mu\one+tW$; the $\ell^1$ contraction for $\mu(P)>1/(2m)$ & Sections~\ref{sec:upper}, \ref{sec:lower}\\
\ref{sec:binary} & $m=2$: discriminant, second iterate, two-cycle & Theorem~\ref{thm:binary}\\
\ref{sec:upper} & Fisher information bound, uniform mixture bound, contraction & Theorem~\ref{thm:upper}\\
\ref{sec:lower} & a positive operator with a two-cycle for each $\ell<L_m$ & Theorem~\ref{thm:lower}\\
\ref{sec:numerics} & optimization of two-cycles and fixed-point pairs for $m=3,4$ & Table~\ref{tab:search}\\
\bottomrule
\end{tabular}
\end{center}

\section{The uniform part of the coefficients}\label{sec:prelim}

Let $\mu=\mu(P)$. If $\mu=1/m$, then every $P_{ij,k}$ equals $1/m$ and $V$ is constant. Otherwise put $t=1-m\mu>0$ and
\[
Q_{ij,k}=\frac{P_{ij,k}-\mu}{t}.
\]
The numbers $Q_{ij,k}$ again satisfy~\eqref{eq:coeff}. Let $W$ be the operator~\eqref{eq:qso} with coefficients
$Q_{ij,k}$, and write $Q(x,y)$ for the symmetric bilinear map with $Q(x,y)_k=\sum_{i,j}Q_{ij,k}x_iy_j$, so that
$W(x)=Q(x,x)$. Since $\sum_{i,j}x_ix_j=1$ on $S^{m-1}$,
\begin{equation}\label{eq:split}
V(x)=\mu\one+tW(x),\qquad x\in S^{m-1},
\end{equation}
where $\one=(1,\dots,1)$. In particular $(Vx)_k\ge\mu$ for every $x$ and $k$, so $V$ maps $S^{m-1}$ into
$\Delta_\mu=\{x\in S^{m-1}:x_k\ge\mu\text{ for all }k\}$.

\begin{proposition}\label{prop:l1}
$\|Vx-Vy\|_1\le2t\|x-y\|_1$ for $x,y\in S^{m-1}$. Hence $V$ is regular if $\mu(P)>1/(2m)$.
\end{proposition}

\begin{proof}
By symmetry, $W(x)-W(y)=Q(x,x-y)+Q(x-y,y)$. For a probability vector $z$ and any $h\in\R^m$,
$\|Q(z,h)\|_1\le\sum_{i,j,k}Q_{ij,k}z_i|h_j|=\|h\|_1$. Hence $\|Vx-Vy\|_1\le2t\|x-y\|_1$, and $2t<1$ when
$\mu>1/(2m)$; the Banach contraction principle on the compact set $S^{m-1}$ gives a fixed point $p$ with $V^nx\to p$ for
every $x$.
\end{proof}

\section{Two types}\label{sec:binary}

Let $m=2$, write $x$ for the first coordinate of a point of $S^1$, and put $a=P_{11,1}$, $b=P_{12,1}=P_{21,1}$ and
$c=P_{22,1}$. The coefficients with $k=2$ are $1-a$, $1-b$ and $1-c$, so $\mu(P)\ge\delta$ means
$a,b,c\in[\delta,1-\delta]$. The first coordinate of $Vx$ is
\begin{equation}\label{eq:f}
f(x)=ax^2+2bx(1-x)+c(1-x)^2=Ax^2+Bx+C,\qquad A=a-2b+c,\quad B=2(b-c),\quad C=c .
\end{equation}
Let $\delta_*=(3-\sqrt7)/4$, the smaller root of $8\delta^2-12\delta+1=0$.

\begin{lemma}\label{lem:disc}
Let $D=(B-1)^2-4AC$ be the discriminant of $f(x)-x$. Then $D=(2b-1)^2+4c(1-a)$. If $a,b,c\in[\delta,1-\delta]$ with
$0<\delta\le1/2$, then $D\le5-12\delta+8\delta^2$, with equality for $a=b=\delta$, $c=1-\delta$. In particular $D\le4$
when $\delta\ge\delta_*$.
\end{lemma}

\begin{proof}
The identity is a direct expansion. Then $(2b-1)^2\le(1-2\delta)^2$ and $c(1-a)\le(1-\delta)^2$, both with equality at
the stated point, and $(1-2\delta)^2+4(1-\delta)^2=5-12\delta+8\delta^2$. This is at most $4$ exactly when
$8\delta^2-12\delta+1\le0$, which for $\delta\le1/2$ means $\delta\ge\delta_*$.
\end{proof}

\begin{lemma}\label{lem:second}
$f(f(x))-x=(f(x)-x)H(x)$, where
\[
H(x)=A^2x^2+A(B+1)x+AC+B+1=\Bigl(Ax+\frac{B+1}{2}\Bigr)^2+\frac{4-D}{4}.
\]
\end{lemma}

\begin{proof}
Both identities are direct expansions.
\end{proof}

\begin{proof}[Proof of Theorem~\ref{thm:binary}]
Assume $a,b,c\in[\delta,1-\delta]$ with $\delta\ge\delta_*$. Since $f(0)=c>0$ and $f(1)=a<1$, the function $f(x)-x$ has
exactly one zero $p$ in $[0,1]$; it lies in $(0,1)$, and $f(x)>x$ for $x<p$, $f(x)<x$ for $x>p$. (If $A\ne0$ this holds
because a quadratic polynomial that changes sign on $[0,1]$ has exactly one zero there; if $A=0$ the function is
affine.)

We claim that
\begin{equation}\label{eq:f2}
f(f(x))>x\ \text{ for } x<p,\qquad f(f(x))<x\ \text{ for } x>p .
\end{equation}
If $A=0$, then $f(x)=Bx+C$ with $B=2(b-c)=a-c$, so $|B|\le1-2\delta<1$, and $f(f(x))-x=(1+B)(f(x)-x)$ has the sign of $f(x)-x$. Let $A\ne0$. By
Lemma~\ref{lem:disc}, $D\le4$. If $D<4$, then $H>0$ and~\eqref{eq:f2} follows from Lemma~\ref{lem:second}. If $D=4$,
then $H$ vanishes only at $r=-(B+1)/(2A)$, and substituting $r$ into $f(x)-x=Ax^2+(B-1)x+C$ gives $(4-D)/(4A)=0$.
So $r$ is a zero of $f(x)-x$, hence $r=p$ or $r\notin[0,1]$, and $H>0$ on $[0,1]\setminus\{p\}$; again~\eqref{eq:f2}
follows.

Let $x\in[0,1]$ and let $\Omega$ be the set of limit points of $(f^n(x))_{n\ge0}$. It is compact and non-empty, and
$f(\Omega)=\Omega$: a limit point of $(f^{n_j}(x))$ is the image of a limit point of $(f^{n_j-1}(x))$, and the image of
a limit point is a limit point. Let $a_0=\min\Omega$ and $b_0=\max\Omega$. Suppose $a_0<b_0$. If $\Omega\subset[p,1]$,
then $b_0>p$ and $b_0=f(y)$ with $y\in\Omega$; but $f(p)=p<b_0$ and $f(y)<y\le b_0$ for $y>p$. The case
$\Omega\subset[0,p]$ is symmetric, so $a_0<p<b_0$. Choose $u,v\in\Omega$ with $f(u)=b_0$ and $f(v)=a_0$; the same
argument gives $u<p<v$. Since $f(u)=b_0\ge v>p=f(p)$, there is $w\in[u,p)$ with $f(w)=v$. Then
$f(f(w))=a_0\le u\le w$, which contradicts~\eqref{eq:f2}. Hence $a_0=b_0$, so $f^n(x)$ converges, and its limit is a
fixed point of $f$, that is, $p$. Thus $V$ is regular whenever $\mu(P)\ge\delta_*$, and $\lambda_2\le\delta_*$.

Now let $0<\delta<\delta_*$ and $a=b=\delta$, $c=1-\delta$, so that $\mu(P)=\delta$ and, with $s=1-2\delta$,
\begin{equation}\label{eq:fdelta}
f(x)=\delta+s(1-x)^2 .
\end{equation}
Put $R=4s(1-\delta)-3$ and
\[
x_\pm=\frac{2s-1\pm\sqrt R}{2s}.
\]
Here $R=D-4$, which is positive by Lemma~\ref{lem:disc} because $5-12\delta+8\delta^2>4$, and $R<1$. Moreover $2s-1>0$
and $(2s-1)^2-R=4\bigl((1-s)^2+s\delta\bigr)>0$, so $0<x_-<x_+<1$. Substitution in~\eqref{eq:fdelta} gives
$f(x_+)=x_-$ and $f(x_-)=x_+$; equivalently, $x_\pm$ are the zeros of the quadratic factor $H$ of
Lemma~\ref{lem:second}. So $V$ has a two-cycle and is not regular. Given $\lambda<\delta_*$, choose
$\delta\in(\max(\lambda,0),\delta_*)$; all coefficients of this operator, namely $\delta$ and $1-\delta$, exceed
$\lambda$. Hence $\lambda_2=\delta_*$.
\end{proof}

\begin{remark}\label{rem:binaryupper}
For $m=2$, Theorem~\ref{thm:upper} gives the regularity of $V$ for $\mu(P)>C_2=(2-\sqrt2)/4=0.146\dots$, a weaker
statement than Theorem~\ref{thm:binary}, by a different argument. The construction of Section~\ref{sec:lower} needs a
third type.
\end{remark}

\section{A contraction in the Fisher--Rao metric}\label{sec:upper}

For $x$ in the interior of $S^{m-1}$ and a tangent vector $h\in\R^m$ with $\sum_ih_i=0$, let
\[
|h|_x^2=\sum_{i=1}^m\frac{h_i^2}{x_i}.
\]
This is the Fisher--Rao metric. The map $x\mapsto2\sqrt x=(2\sqrt{x_1},\dots,2\sqrt{x_m})$ carries it to the round metric
on the sphere of radius $2$, so $S^{m-1}$ is isometric to the closed positive orthant of that sphere, with the distance
\begin{equation}\label{eq:dist}
d(x,y)=2\arccos\Bigl(\sum_{i=1}^m\sqrt{x_iy_i}\Bigr).
\end{equation}
The distance $d$ is continuous on $S^{m-1}\times S^{m-1}$ and induces the usual topology, so $(S^{m-1},d)$ is compact
and complete. For $x,y$ in the interior, the shorter great-circle arc between $2\sqrt x$ and $2\sqrt y$ stays in the
open orthant, and its length is $d(x,y)$.

\begin{lemma}\label{lem:fisher}
Let $W$ be a quadratic stochastic operator with coefficients $Q_{ij,k}$, let $x$ be in the interior of $S^{m-1}$ and let
$\sum_ih_i=0$. Then
\[
\sum_{k:\,W_k(x)>0}\frac{(DW(x)h)_k^2}{W_k(x)}\le2|h|_x^2,
\]
and $(DW(x)h)_k=0$ whenever $W_k(x)=0$.
\end{lemma}

\begin{proof}
Put $u_i=h_i/x_i$. By symmetry of $Q$, $(DW(x)h)_k=2Q(x,h)_k=\sum_{i,j}Q_{ij,k}x_ix_j(u_i+u_j)$. If $W_k(x)=0$, then
$Q_{ij,k}=0$ for all $i,j$, because every $x_ix_j$ is positive, and so $(DW(x)h)_k=0$. Otherwise the Cauchy--Schwarz
inequality with the weights $Q_{ij,k}x_ix_j$, which sum to $W_k(x)$, gives
\[
(DW(x)h)_k^2\le W_k(x)\sum_{i,j}Q_{ij,k}x_ix_j(u_i+u_j)^2 .
\]
Divide by $W_k(x)$ and sum over $k$; since $\sum_kQ_{ij,k}=1$,
\[
\sum_k\frac{(DW(x)h)_k^2}{W_k(x)}\le\sum_{i,j}x_ix_j(u_i+u_j)^2=2\sum_ix_iu_i^2+2\Bigl(\sum_ix_iu_i\Bigr)^2=2|h|_x^2,
\]
because $\sum_ix_iu_i=\sum_ih_i=0$.
\end{proof}

\begin{lemma}\label{lem:mix}
Let $q$ be a probability vector with positive entries, let $v\in\R^m$ with $\sum_kv_k=0$, and let $\mu>0$ and
$t\ge0$. Then
\[
\sum_k\frac{v_k^2}{\mu+tq_k}\le\frac{1}{t+2\mu}\sum_k\frac{v_k^2}{q_k}.
\]
The constant is sharp. It is attained for $m=2$ at $q=(1/2,1/2)$, and for $m>2$ it is approached as $q$ tends to
$(1/2,1/2,0,\dots,0)$.
\end{lemma}

\begin{proof}
The right side minus the left side is $\frac{\mu}{t+2\mu}\sum_kv_k^2\frac{1-2q_k}{q_k(\mu+tq_k)}$, so it suffices to
show
\begin{equation}\label{eq:mix}
\sum_kv_k^2\frac{1-2q_k}{q_k(\mu+tq_k)}\ge0 .
\end{equation}
If $q_k\le1/2$ for every $k$, each term is non-negative. Otherwise exactly one entry exceeds $1/2$, say $q_1=s>1/2$, and
the others are positive and sum to $r=1-s<1/2$. For $k\ge2$ put $a_k=(1-2q_k)/(q_k(\mu+tq_k))>0$. The function
$\varphi(q)=q(\mu+tq)/(1-2q)$ is non-negative, increasing and convex on $[0,1/2)$ and vanishes at $0$, so it is
superadditive there, and
\[
\sum_{k\ge2}\frac1{a_k}=\sum_{k\ge2}\varphi(q_k)\le\varphi(r)=\frac{r(\mu+tr)}{2s-1}\le\frac{s(\mu+ts)}{2s-1}.
\]
By the Cauchy--Schwarz inequality and $\sum_{k\ge2}v_k=-v_1$,
\[
\sum_{k\ge2}a_kv_k^2\ge\frac{v_1^2}{\sum_{k\ge2}a_k^{-1}}\ge v_1^2\,\frac{2s-1}{s(\mu+ts)},
\]
which is the negative of the term $k=1$ in~\eqref{eq:mix}. For $q=(1/2,1/2,0,\dots,0)$ restricted to its support and
$v=(1,-1,0,\dots,0)$, both sides equal $4/(t+2\mu)$.
\end{proof}

\begin{proof}[Proof of Theorem~\ref{thm:upper}]
Let $\mu=\mu(P)>C_m$. If $\mu=1/m$, $V$ is constant. Otherwise write $V=\mu\one+tW$ as in~\eqref{eq:split}, with
$t=1-m\mu$. Let $x$ be in the interior of $S^{m-1}$ and $\sum_ih_i=0$. Then $DV(x)h=t\,DW(x)h$ and $V_k(x)=\mu+tW_k(x)$.
Apply Lemma~\ref{lem:mix} with $q=W(x)$ and $v=DW(x)h$, after deleting the coordinates with $W_k(x)=0$ (there $v_k=0$
by Lemma~\ref{lem:fisher}, and the remaining entries of $q$ still sum to $1$ and those of $v$ to $0$), and then
Lemma~\ref{lem:fisher}:
\begin{equation}\label{eq:K}
|DV(x)h|_{V(x)}^2=t^2\sum_k\frac{(DW(x)h)_k^2}{\mu+tW_k(x)}\le\frac{t^2}{t+2\mu}\sum_{k:\,W_k(x)>0}\frac{(DW(x)h)_k^2}{W_k(x)}\le K^2|h|_x^2,
\end{equation}
where
\[
K^2=\frac{2t^2}{t+2\mu}=\frac{2(1-m\mu)^2}{1-(m-2)\mu}.
\]
Let $x,y$ be in the interior and let $\gamma$ be the curve in $S^{m-1}$ whose image under $x\mapsto2\sqrt x$ is the
shorter great-circle arc from $2\sqrt x$ to $2\sqrt y$. By~\eqref{eq:K} the curve $V\circ\gamma$, which joins $Vx$ to
$Vy$, has Fisher--Rao length at most $K\,d(x,y)$, and $d(Vx,Vy)$ is at most that length because great-circle arcs are
shortest on the sphere. So $d(Vx,Vy)\le K\,d(x,y)$ on the interior, and by continuity on all of $S^{m-1}$.

Finally, $K<1$ is equivalent to $2(1-m\mu)^2<1-(m-2)\mu$, that is, to
\[
2m^2\mu^2-(3m+2)\mu+1<0 .
\]
The roots of this quadratic are $\bigl(3m+2\mp\sqrt{m^2+12m+4}\bigr)/(4m^2)$. The smaller one is $C_m$, and the larger
one is at least $1/m$ because $\sqrt{m^2+12m+4}\ge m-2$. Hence $K<1$ for $C_m<\mu<1/m$, $V$ is a contraction of the
complete metric space $(S^{m-1},d)$, and the Banach contraction principle shows that $V$ is regular. If all coefficients
exceed $C_m$, then so does $\mu(P)$; thus $\lambda_m\le C_m$. At $\mu=1/(2m)$ one has $K^2=m/(m+2)<1$, and
$1/(2m)<1/m$, so $C_m<1/(2m)$.
\end{proof}

\begin{remark}\label{rem:l1}
The set $\Delta_\mu$ also gives a sharper $\ell^1$ bound than Proposition~\ref{prop:l1}. For $x\in\Delta_\mu$ let
$B_{ki}=\sum_jQ_{ij,k}x_j$. The columns $i$ and $l$ of the stochastic matrix $B$ have overlap
$\sum_k\min(B_{ki},B_{kl})\ge\sum_kQ_{il,k}\min(x_l,x_i)\ge\mu$, so the Dobrushin coefficient of $B$ is at most $1-\mu$.
Since $DV(x)=2tB$ on the tangent space, integration along segments in the convex set $\Delta_\mu$ shows that $V$ contracts $\Delta_\mu$ in the
$\ell^1$ norm with factor $2(1-m\mu)(1-\mu)$, which is less than $1$ for $\mu>1/(m+1+\sqrt{m^2+1})$. This threshold
lies between $C_m$ and $1/(2m)$.
\end{remark}

\section{Positive operators with a two-cycle}\label{sec:lower}

Let $m\ge3$. Divide the types into the classes $\mathcal A=\{1\}$, $\mathcal C=\{3\}$ and $\mathcal B=\{2,4,5,\dots,m\}$,
and let $Q$ be the deterministic heredity rule in which the offspring of parents from the classes
\[
\mathcal A\mathcal A\text{ or }\mathcal A\mathcal B\ \text{ is of type }3,\qquad
\mathcal C\mathcal C\text{ or }\mathcal C\mathcal B\ \text{ is of type }1,\qquad
\mathcal A\mathcal C\text{ or }\mathcal B\mathcal B\ \text{ is of type }2 .
\]
That is, $Q_{ij,k}=1$ if $k$ is the type assigned to the classes of $i$ and $j$, and $Q_{ij,k}=0$ otherwise. For
$0<\ell<1/m$ put $\tau=1-m\ell$ and
\begin{equation}\label{eq:P}
P_{ij,k}=\ell+\tau Q_{ij,k}.
\end{equation}
These coefficients satisfy~\eqref{eq:coeff}, take only the values $\ell$ and $1-(m-1)\ell$, and have $\mu(P)=\ell$.

\begin{proof}[Proof of Theorem~\ref{thm:lower}]
Write $a=x_1$, $c=x_3$ and $z=\sum_{i\in\mathcal B}x_i=1-a-c$. By~\eqref{eq:split} and the rule above, the operator
$V$ defined by~\eqref{eq:P} is
\begin{equation}\label{eq:map}
(Vx)_1=\ell+\tau(c^2+2cz),\qquad (Vx)_3=\ell+\tau(a^2+2az),\qquad (Vx)_k=\ell\ \ (k\ge4),
\end{equation}
and $(Vx)_2$ is determined by $\sum_k(Vx)_k=1$. Let $\ell\in(0,L_m)$, so that
\begin{equation}\label{eq:tau}
\tau>\frac12\Bigl(1+\frac1{\sqrt{m+1}}\Bigr).
\end{equation}
Put
\[
z_0=\frac1\tau-1=\frac{m\ell}{\tau},\qquad \sigma=2-\frac1\tau,\qquad
\Lambda=\sigma^2-\frac{4\ell}{\tau}=\frac{(m+1)(2\tau-1)^2-1}{m\tau^2},
\]
where the last expression follows from $\ell=(1-\tau)/m$. By~\eqref{eq:tau}, $\sigma>0$ and $\Lambda>0$, so the numbers
$u=(\sigma+\sqrt\Lambda)/2$ and $v=(\sigma-\sqrt\Lambda)/2$ are real and distinct, with $u+v=\sigma$ and $uv=\ell/\tau>0$;
hence $u,v>0$. Since $z_0/\ell=m/\tau>m-3$, the points
\[
X=\bigl(u,\ z_0-(m-3)\ell,\ v,\ \ell,\dots,\ell\bigr),\qquad Y=\bigl(v,\ z_0-(m-3)\ell,\ u,\ \ell,\dots,\ell\bigr)
\]
have positive coordinates, and their coordinates sum to $u+v+z_0=1$; so $X\ne Y$ lie in $S^{m-1}$, and for both of them
$z=z_0$. Using $\tau(1+z_0)=1$, $u+v+2z_0=1/\tau$ and $\ell=\tau uv$, \eqref{eq:map} gives at $X$
\[
(VX)_1=\tau uv+\tau v(v+2z_0)=\tau v(u+v+2z_0)=v,
\]
and in the same way $(VX)_3=u$. The coordinates $k\ge4$ equal $\ell$, and then $(VX)_2=z_0-(m-3)\ell$. So $VX=Y$, and by
the symmetry between the types $1$ and $3$ in the rule, $VY=X$. The operator $V$ has the two-cycle $\{X,Y\}$ and is not
regular, and $\mu(P)=\ell$. Given $\lambda<L_m$, choose $\ell\in(\max(\lambda,0),L_m)$; all coefficients
of~\eqref{eq:P} exceed $\lambda$. Hence $\lambda_m\ge L_m$.
\end{proof}

\begin{remark}\label{rem:bifurcation}
The two-cycle is born from the symmetric fixed point at $\ell=L_m$. In the coordinates $e=a-c$ and $z$ the map
\eqref{eq:map} reads $e'=-\tau(1+z)e$ and $z'=(m-2)\ell+\frac\tau2(3z^2-2z+1-e^2)$, and the derivative at the fixed
point with $e=0$ has the eigenvalue $-1$ in the direction of $e$ exactly when $\ell=L_m$. If the offspring types $1$ and
$3$ are exchanged in the rule, the first equation becomes $e'=\tau(1+z)e$, and $X$, $Y$ are two distinct fixed points
instead of a two-cycle.
\end{remark}

\begin{remark}\label{rem:examples}
For some rational $\ell$ the cycle has explicit coordinates. For $m=3$ and $\ell=2/25$,
$X=\bigl((13+\sqrt{17})/38,\ 6/19,\ (13-\sqrt{17})/38\bigr)$ and the coefficients are $2/25$ and $21/25$. For $m=4$ and
$\ell=11/160$, $X=\bigl((18+\sqrt5)/58,\ 1441/4640,\ (18-\sqrt5)/58,\ 11/160\bigr)$. For $m=5$ and $\ell=7/120$,
$X=\bigl((25+\sqrt{30})/85,\ 301/1020,\ (25-\sqrt{30})/85,\ 7/120,\ 7/120\bigr)$. In these three examples the cycle
attracts a neighbourhood of itself: the derivative of $V^2$ at $X$, restricted to the tangent space of $S^{m-1}$, has
spectral radius $0.947$, $0.993$ and $0.980$ respectively, computed from its exact characteristic polynomial.
\end{remark}

\section{Numerical searches}\label{sec:numerics}

Table~\ref{tab:search} gives, for $m=3,4$, the largest least coefficient reached by numerical optimization of an
operator with a two-cycle or with two distinct fixed points. For given points $X\ne Y$ of $S^{m-1}$, the largest $\mu(P)$
subject to $VX=Y$ and $VY=X$ (or to $VX=X$ and $VY=Y$) is the value of a linear program in the coefficients. The points
are optimized by differential evolution under a lower bound on the separation $|X_1-Y_1|$, which keeps a fixed point
from counting as a cycle, and the best results are then refined in all variables together. The values approach $L_m$
from below as the separation decreases, and none exceeds $L_m$. In random samples of $4{,}320$ operators with least
coefficient between $L_m+10^{-4}$ and $C_m+10^{-4}$, the trajectories from the vertices and from four random points
converge to a common limit. These computations do not exclude cycles of longer period, non-periodic invariant sets or a
different family with least coefficient in $(L_m,C_m]$.

\begin{table}[tbp]
\centering
\begin{tabular}{rcccc}
\toprule
$m$ & separation & largest $\mu(P)$ found & $L_m$ & $C_m$\\
\midrule
3 & $0.04$ & $0.083221$ & $0.083333$ & $0.111111$\\
3 & $0.008$ & $0.083329$ & $0.083333$ & $0.111111$\\
4 & $0.04$ & $0.069005$ & $0.069098$ & $0.089903$\\
4 & $0.008$ & $0.069095$ & $0.069098$ & $0.089903$\\
\bottomrule
\end{tabular}
\caption{Largest least coefficient of an operator with a two-cycle or two distinct fixed points reached by the
optimization of Section~\ref{sec:numerics}, for two lower bounds on the separation $|X_1-Y_1|$.}\label{tab:search}
\end{table}

\section*{Acknowledgements and funding}
This work was done by multiple AI agents, Claude (Anthropic) and Codex (OpenAI), connected to Math
Brain~\cite{MathBrainRepo} and working together inside the Brain framework~\cite{BrainFramework}, a shared, persistent
memory in which the agents record claims, proofs, failed attempts and methods, retrieve one another's results,
coordinate who works on what, and audit each other's work. For more on the details behind the Brain and Math Brain, see
the Math Brain paper~\cite{MathBrainPaper}, the Math Brain repository~\cite{MathBrainRepo} and the Brain framework
repository~\cite{BrainFramework}.

The research presented in this paper was supported by the European Research Council (ERC) under the European Union's
Horizon Europe research and innovation programme (grant agreement No.~101041711), by the Simons Foundation as part of
the Collaboration on the Mathematical and Scientific Foundations of Deep Learning, by Heights Labs, by Convex Nexus
Capital, by the Israel Science Foundation (grant number 2258/19), and by the Israel Science Foundation (ISF Grant
4101/25). Computational resources were provided in part by Google Cloud through its education credits program (credit
530550103).

\section*{Declaration of competing interest}
The author holds equity in Heights Labs and in Convex Nexus Capital and has a management role at Convex Nexus Capital.
These entities are included in the funding declaration above.

\section*{Declaration of generative AI and AI-assisted technologies}
A Codex agent (OpenAI) proved Theorems~\ref{thm:binary}, \ref{thm:upper} and~\ref{thm:lower}; a second Codex agent
checked the proofs as a referee, gave the $\ell^1$ bound of Remark~\ref{rem:l1} its proof and ran the searches of
Section~\ref{sec:numerics}; Claude agents (Anthropic) checked the steps and prepared the manuscript. The problem and the bound $1/(2m)$ are due
to Rasul N. Ganikhodzhaev and A. T. Sarymsakov~\cite{GS88}, Ganikhodzhaev and Utkir A. Rozikov~\cite{GR09} and
Ganikhodzhaev, Farrukh Mukhamedov and Rozikov~\cite{GMR11}; the binary discriminant criterion to Yu. I.
Lyubich~\cite{Lyubich74}, as recalled by Karim, Hamzah and Nasir Ganikhodjaev~\cite{KHG22}; positive operators with
several fixed points to Krapivin and Lyubich~\cite{KL76}; the ergodicity coefficient to Dobrushin~\cite{Dobrushin56} and
the optimal $\ell^1$ contraction to Fasino and Tudisco~\cite{FT20}; the doubly stochastic case to M.~Kh. Saburov and
Kh.~Kh. Saburov~\cite{SS25}. New here are the value of $\lambda_2$ with its boundary case, the Fisher--Rao contraction
behind Theorem~\ref{thm:upper} and the two-cycles behind Theorem~\ref{thm:lower}. The author directed
the work and is responsible for the content of the paper and for decisions concerning its dissemination.

\section*{Code availability}
Exact checks of the identities in Sections~\ref{sec:binary}--\ref{sec:lower}, of the two-cycles for $m=3,\dots,8$ and of
the examples of Remark~\ref{rem:examples}, and the searches of Section~\ref{sec:numerics}, are in the repository
\begin{center}\url{https://github.com/yspennstate/qso-regularity-threshold/tree/v1.1}.\end{center}

\appendix
\section{Notation}\label{app:notation}

\begin{center}
\small
\begin{tabular}{lll}
\toprule
Symbol & Meaning & First used\\
\midrule
$S^{m-1}$ & the simplex of probability vectors in $\R^m$ & Section~\ref{sec:intro}\\
$P_{ij,k}$ & heredity coefficients of $V$ & \eqref{eq:qso}\\
$\mu(P)$, $\mu$ & the least heredity coefficient & Section~\ref{sec:intro}\\
$\lambda_m$ & the regularity threshold & \eqref{eq:lambda}\\
$C_m$, $L_m$ & the upper and lower bounds for $\lambda_m$ & Theorems~\ref{thm:upper}, \ref{thm:lower}\\
$t$ & $1-m\mu$ & Section~\ref{sec:prelim}\\
$Q_{ij,k}$, $W$, $Q(x,y)$ & $P$, $V$ and the bilinear map without the uniform part & \eqref{eq:split}\\
$\one$ & the vector $(1,\dots,1)$ & \eqref{eq:split}\\
$\Delta_\mu$ & points of $S^{m-1}$ with all coordinates $\ge\mu$ & Section~\ref{sec:prelim}\\
$a,b,c$; $A,B,C$ & coefficients of the binary map $f$ & \eqref{eq:f}\\
$D$ & the discriminant of $f(x)-x$ & Lemma~\ref{lem:disc}\\
$\delta_*$ & $(3-\sqrt7)/4$ & Section~\ref{sec:binary}\\
$H$ & the quadratic factor of $f(f(x))-x$ & Lemma~\ref{lem:second}\\
$|h|_x$, $d$ & the Fisher--Rao norm and distance & \eqref{eq:dist}\\
$K$ & the Lipschitz constant of $V$ for $d$ & \eqref{eq:K}\\
$\mathcal A,\mathcal B,\mathcal C$ & classes of types & Section~\ref{sec:lower}\\
$\ell$, $\tau$ & the least coefficient and $1-m\ell$ in~\eqref{eq:P} & \eqref{eq:P}\\
$z_0,\sigma,\Lambda,u,v$ & parameters of the two-cycle $X,Y$ & Section~\ref{sec:lower}\\
\bottomrule
\end{tabular}
\end{center}

\begin{thebibliography}{99}
\bibitem{FT20} D.~Fasino and F.~Tudisco, Ergodicity coefficients for higher-order stochastic processes, \emph{SIAM J.\
Math.\ Data Sci.} \textbf{2} (2020), no.~3, 740--769. \doi{10.1137/19M1285214}
\bibitem{GMR11} R.~Ganikhodzhaev, F.~Mukhamedov and U.~Rozikov, Quadratic stochastic operators and processes: results
and open problems, \emph{Infin.\ Dimens.\ Anal.\ Quantum Probab.\ Relat.\ Top.} \textbf{14} (2011), no.~2, 279--335.
\doi{10.1142/S0219025711004365}
\bibitem{GR09} R.~N.~Ganikhodzhaev and U.~A.~Rozikov, Quadratic stochastic operators: results and open problems,
preprint (2009), \href{https://arxiv.org/abs/0902.4207}{arXiv:0902.4207}.
\bibitem{GS88} R.~N.~Ganikhodzhaev and A.~T.~Sarymsakov, A simple criterion for regularity of quadratic stochastic
operators, \emph{Dokl.\ Akad.\ Nauk UzSSR} (1988), no.~11, 5--6 (in Russian).
\bibitem{KHG22} S.~N.~Karim, N.~Z.~A.~Hamzah and N.~Ganikhodjaev, On the dynamics of geometric quadratic stochastic
operator generated by 2-partition on countable state space, \emph{Malays.\ J.\ Math.\ Sci.} \textbf{16} (2022), no.~4,
727--737. \doi{10.47836/mjms.16.4.06}
\bibitem{Dobrushin56} R.~L.~Dobrushin, Central limit theorem for nonstationary Markov chains. I, \emph{Theory Probab.\
Appl.} \textbf{1} (1956), no.~1, 65--80. \doi{10.1137/1101006}
\bibitem{KL76} A.~A.~Krapivin and Yu.~I.~Lyubich, Estimate of the dimension of the set of fixed points of a class of
quadratic operators, \emph{Sib.\ Math.\ J.} \textbf{17} (1976), no.~4, 704--707. \doi{10.1007/BF00971683}
\bibitem{KL77} A.~A.~Krapivin and Yu.~I.~Lyubich, Estimates of Lipschitz constants for polynomial operators in a
simplex, \emph{Soviet Math.\ Dokl.} \textbf{18} (1977), 718--721.
\bibitem{Lyubich74} Yu.~I.~Lyubich, Iterations of quadratic mappings, in \emph{Mathematical Economics and Functional
Analysis}, Nauka, Moscow, 1974 (in Russian).
\bibitem{S26} M.~Saburov, Monte-Carlo methods for stochastic hyper-matrices and global asymptotic stability of
polynomial stochastic operators, \emph{Funct.\ Anal.\ Appl.} \textbf{60} (2026), no.~3, 312--318.
\doi{10.1134/S1234567826030109}
\bibitem{SS25} M.~Kh.~Saburov and Kh.~Kh.~Saburov, The nonlinear Perron--Frobenius stability problem for cubic stochastic
matrices, \emph{Russian Math.\ Surveys} \textbf{80} (2025), no.~5, 931--933. \doi{10.4213/rm10256e}
\bibitem{MathBrainRepo} Y.~Shmalo, \emph{Math Brain: public snapshot, code and tests}, 2026,
\url{https://github.com/yspennstate/math-brain-public/tree/8c9a6bd478c4ec4082347151b7ae16db50c917dd}.
\bibitem{MathBrainPaper} Y.~Shmalo, \emph{Math Brain: shared, auditable memory for multi-agent mathematical research},
October 2026, \url{https://github.com/yspennstate/math-brain-public/tree/main/papers/math-brain-2026-10-24},
\doi{10.5281/zenodo.24150837}.
\bibitem{BrainFramework} Y.~Shmalo, \emph{Brain Framework}, public software repository, 2026,
\url{https://github.com/yspennstate/brain-framework-public/tree/5629e644cdd510173ae97e9aa92642dabd22dc57}.
\end{thebibliography}
\end{document}
